\section{General Polynomial Potentials}\label{sec:general}
% ============================================================
% ============================================================

The previous sections treated the quartic double-well $V(x)=ax^4-bx^2$ as the
canonical example. This section develops the complete analytical machinery for
an \emph{arbitrary} polynomial potential of degree $N$, an arbitrary number of
basis functions $M$, and user-defined initial wavepackets. The derivations here
underpins the \texttt{QuTu} general-case solver.

% ============================================================
\subsection{General polynomial potential}\label{sec:gen_pot}
% ============================================================

A general real-valued polynomial potential of even degree $N$ (with $N \geq 2$)
is written
\begin{equation}\label{eq:gen_V}
  V(x) = \sum_{k=0}^{N} c_k\,x^k, \qquad c_N > 0,
\end{equation}
with real coefficients $\{c_k\}$.  The requirement $c_N > 0$ ensures
confinement ($V(x)\to+\infty$ as $|x|\to\infty$).  The quartic double-well
of Section~\ref{sec:potential} corresponds to $N=4$, $c_4=a$, $c_2=-b$, and
all other coefficients zero.

\begin{remark}
Odd-degree leading terms are excluded because $V(x)\to-\infty$ along one
direction, yielding an unbounded-below Hamiltonian with no ground state.
Odd \emph{interior} coefficients ($c_1,c_3,\ldots$) are perfectly allowed;
they break the left--right symmetry of the potential.
\end{remark}

% ============================================================
\subsection{Hamiltonian matrix in the harmonic-oscillator basis}\label{sec:gen_H}
% ============================================================

Using the harmonic oscillator (HO) basis $\{\phi_n\}_{n=0}^{M-1}$ introduced
in Section~\ref{sec:HO}, the Hamiltonian matrix elements decompose as
\begin{equation}\label{eq:gen_H_elem}
  H_{mn} = T_{mn} + V_{mn},
\end{equation}
where $T_{mn}=\bra{m}\hat T\ket{n}$ is the kinetic energy matrix element
(derived in Section~\ref{sec:HO}) and
\begin{equation}\label{eq:gen_V_elem}
  V_{mn} = \bra{m}V(\hat x)\ket{n}
         = \sum_{k=0}^{N} c_k\, X^{(k)}_{mn},
  \qquad
  X^{(k)}_{mn} \equiv \bra{m}\hat x^k\ket{n}.
\end{equation}
The problem therefore reduces to computing all power-matrix elements
$X^{(k)}_{mn}$ for $k=0,1,\ldots,N$.

% ============================================================
\subsection{Power-matrix elements via the master recursion}\label{sec:master_rec}
% ============================================================

\subsubsection{Ladder-operator form of \texorpdfstring{$\hat x$}{x}}

Introduce the characteristic length
\begin{equation}\label{eq:ell}
  \ell = \frac{1}{\sqrt{2\mu\alpha}},
\end{equation}
where $\mu$ is the mass and $\alpha$ the HO width parameter.  Then
\begin{equation}\label{eq:x_ladder}
  \hat x = \ell\,(\hat a + \hat a^\dagger),
\end{equation}
with the standard actions $\hat a\ket{n}=\sqrt{n}\,\ket{n-1}$ and
$\hat a^\dagger\ket{n}=\sqrt{n+1}\,\ket{n+1}$.

\subsubsection{Master recursion theorem}

\begin{result}[Master recursion for $X^{(k)}_{mn}$]\label{res:master_rec}
  For $k \geq 1$ and arbitrary non-negative integers $m,n$,
  \begin{equation}\label{eq:master_rec}
    X^{(k)}_{mn}
    = \ell\Bigl[\sqrt{n}\,X^{(k-1)}_{m,n-1}
               + \sqrt{n+1}\,X^{(k-1)}_{m,n+1}\Bigr],
  \end{equation}
  with the initial condition $X^{(0)}_{mn} = \delta_{mn}$.
\end{result}

\begin{proof}
Insert $\hat x = \ell(\hat a+\hat a^\dagger)$ and use linearity:
\[
  X^{(k)}_{mn}
  = \bra{m}\hat x^k\ket{n}
  = \bra{m}\hat x^{k-1}\hat x\ket{n}
  = \ell\,\bra{m}\hat x^{k-1}\bigl(\hat a+\hat a^\dagger\bigr)\ket{n}.
\]
Acting with the ladder operators on $\ket{n}$ gives
\[
  X^{(k)}_{mn}
  = \ell\Bigl[\sqrt{n}\,\bra{m}\hat x^{k-1}\ket{n-1}
             + \sqrt{n+1}\,\bra{m}\hat x^{k-1}\ket{n+1}\Bigr]
  = \ell\Bigl[\sqrt{n}\,X^{(k-1)}_{m,n-1}
             + \sqrt{n+1}\,X^{(k-1)}_{m,n+1}\Bigr]. \qed
\]
\end{proof}

Equation~\eqref{eq:master_rec} allows one to compute the entire matrix
$X^{(k)}$ for any $k$ by simple matrix--vector operations; no additional
integrals are required.

% ============================================================
\subsection{Selection rules and bandwidth}\label{sec:selection}
% ============================================================

\begin{result}[Selection rule]\label{res:selection}
  $X^{(k)}_{mn} = 0$ unless the following two conditions both hold:
  \begin{enumerate}
    \item \textbf{Parity:} $(m-n) \equiv k \pmod{2}$.
    \item \textbf{Reach:} $|m-n| \leq k$.
  \end{enumerate}
\end{result}

\begin{proof}
By induction on $k$.  The base case $k=0$ gives $\delta_{mn}$, satisfying
both rules ($(m-n) \equiv 0 \pmod 2$ and $|m-n|=0\leq 0$).

Assume the rules hold for $k-1$.  A non-zero contribution from
$X^{(k-1)}_{m,n\pm 1}$ requires $(m-(n\pm1))\equiv k-1\pmod 2$ and
$|m-(n\pm 1)|\leq k-1$.  The shift $n\to n\pm 1$ changes the parity of
$m-n$ by one unit, so $m-n\equiv k\pmod 2$; and $|m-n|\leq|m-(n\pm1)|+1\leq
k$.  Hence the rules propagate. \qed
\end{proof}

\begin{result}[Bandwidth]\label{res:bandwidth}
  The Hamiltonian matrix $\mathbf{H}$ for the potential~\eqref{eq:gen_V} is
  \emph{banded} with half-bandwidth $N$:
  \begin{equation}
    H_{mn} = 0 \quad \text{if } |m-n| > N.
  \end{equation}
\end{result}

The bandwidth is determined solely by the polynomial degree; the kinetic
energy matrix $\mathbf{T}$ has half-bandwidth 2, which is dominated by $N$
for $N>2$.  For a degree-4 potential, $\mathbf{H}$ has half-bandwidth 4,
consistent with the explicit formulas in Section~\ref{sec:HO}.

% ============================================================
\subsection{Explicit matrix elements for low powers}\label{sec:explicit_Xk}
% ============================================================

Applying the master recursion~\eqref{eq:master_rec} iteratively yields closed
forms for each $k$.  Below, $\ell=(2\mu\alpha)^{-1/2}$ and
$\delta_{mn}\equiv\delta_{m,n}$.

\subsubsection{\texorpdfstring{$k=1$}{k=1}: position matrix}

\begin{equation}\label{eq:X1}
  X^{(1)}_{mn}
  = \ell\bigl(\sqrt{n}\,\delta_{m,n-1} + \sqrt{n+1}\,\delta_{m,n+1}\bigr).
\end{equation}

\subsubsection{\texorpdfstring{$k=2$}{k=2}: quadratic}

\begin{equation}\label{eq:X2}
  X^{(2)}_{mn}
  = \ell^2\Bigl[(2n+1)\,\delta_{mn}
    + \sqrt{n(n-1)}\,\delta_{m,n-2}
    + \sqrt{(n+1)(n+2)}\,\delta_{m,n+2}\Bigr].
\end{equation}

\subsubsection{\texorpdfstring{$k=3$}{k=3}: cubic}

\begin{align}\label{eq:X3}
  X^{(3)}_{mn}
  &= \ell^3\Bigl[3n^{3/2}\,\delta_{m,n-1}
               + 3(n+1)^{3/2}\,\delta_{m,n+1} \nonumber\\
  &\quad +\sqrt{n(n-1)(n-2)}\,\delta_{m,n-3}
    +\sqrt{(n+1)(n+2)(n+3)}\,\delta_{m,n+3}\Bigr].
\end{align}

\subsubsection{\texorpdfstring{$k=4$}{k=4}: quartic}

\begin{align}\label{eq:X4}
  X^{(4)}_{mn}
  &= \ell^4\Bigl[(6n^2+6n+3)\,\delta_{mn}
    + (4n-2)\sqrt{n(n-1)}\,\delta_{m,n-2}
    + (4n+6)\sqrt{(n+1)(n+2)}\,\delta_{m,n+2} \nonumber\\
  &\quad +\sqrt{n(n-1)(n-2)(n-3)}\,\delta_{m,n-4}
    +\sqrt{(n+1)(n+2)(n+3)(n+4)}\,\delta_{m,n+4}\Bigr].
\end{align}

\subsubsection{\texorpdfstring{$k=5$}{k=5}: quintic}

\begin{align}\label{eq:X5}
  X^{(5)}_{mn}
  &= \ell^5\Bigl[(10n^2+10n+5)n^{1/2}\,\delta_{m,n-1}
              + (10n^2+10n+5)(n+1)^{1/2}\,\delta_{m,n+1} \nonumber\\
  &\quad + 5(2n-1)\sqrt{n(n-1)(n-2)}\,\delta_{m,n-3}
    + 5(2n+3)\sqrt{(n+1)(n+2)(n+3)}\,\delta_{m,n+3} \nonumber\\
  &\quad + \sqrt{n(n-1)(n-2)(n-3)(n-4)}\,\delta_{m,n-5} \nonumber\\
  &\quad + \sqrt{(n+1)(n+2)(n+3)(n+4)(n+5)}\,\delta_{m,n+5}\Bigr].
\end{align}

\subsubsection{\texorpdfstring{$k=6$}{k=6}: sextic}

\begin{align}\label{eq:X6}
  X^{(6)}_{mn}
  &= \ell^6\Bigl[(20n^3+30n^2+40n+15)\,\delta_{mn} \nonumber\\
  &\quad + 3(10n^2-10n+4)\sqrt{n(n-1)}\,\delta_{m,n-2}
    + 3(10n^2+10n+4)\sqrt{(n+1)(n+2)}\,\delta_{m,n+2} \nonumber\\
  &\quad + (15n-15)\sqrt{n(n-1)(n-2)(n-3)}\,\delta_{m,n-4} \nonumber\\
  &\quad + (15n+15)\sqrt{(n+1)(n+2)(n+3)(n+4)}\,\delta_{m,n+4} \nonumber\\
  &\quad + \sqrt{n(n-1)(n-2)(n-3)(n-4)(n-5)}\,\delta_{m,n-6} \nonumber\\
  &\quad + \sqrt{(n+1)(n+2)(n+3)(n+4)(n+5)(n+6)}\,\delta_{m,n+6}\Bigr].
\end{align}

\subsubsection{Summary table}

Table~\ref{tab:Xk_bandwidth} summarises the non-zero diagonal offsets and
leading coefficient for each $k$.

\begin{table}[htbp]
  \centering
  \caption{%
    Non-zero off-diagonal levels and leading prefactors for
    $X^{(k)}_{mn}$ ($k=1,\ldots,6$).
    $\ell=(2\mu\alpha)^{-1/2}$.
    ``Offsets'' lists the values of $m-n$ for which $X^{(k)}_{mn}\neq 0$.%
  }
  \label{tab:Xk_bandwidth}
  \ra{1.3}
  \begin{tabular}{ccc}
    \toprule
    $k$ & Non-zero offsets $m-n$ & Prefactor \\
    \midrule
    1 & $\pm 1$ & $\ell$ \\
    2 & $0, \pm 2$ & $\ell^2$ \\
    3 & $\pm 1, \pm 3$ & $\ell^3$ \\
    4 & $0, \pm 2, \pm 4$ & $\ell^4$ \\
    5 & $\pm 1, \pm 3, \pm 5$ & $\ell^5$ \\
    6 & $0, \pm 2, \pm 4, \pm 6$ & $\ell^6$ \\
    \bottomrule
  \end{tabular}
\end{table}

% ============================================================
\subsection{Optimal width parameter \texorpdfstring{$\alpha$}{alpha}}\label{sec:opt_alpha}
% ============================================================

The convergence of the variational expansion depends critically on the choice
of the HO width parameter $\alpha$.  The stationarity condition
$\partial E_0/\partial\alpha = 0$ via the Hellmann--Feynman theorem gives the
implicit equation
\begin{equation}\label{eq:opt_alpha_general}
  \langle\hat T\rangle_\alpha = \frac{1}{2}\sum_{k=0}^{N} k\,c_k\,
  \langle\hat x^k\rangle_\alpha,
\end{equation}
where all expectation values are taken in the ground eigenstate at width
$\alpha$.  This must be solved self-consistently.

\subsubsection{Analytical estimate for a pure monomial}

For the pure monomial $V(x)=c_N x^N$ ($N$ even), the virial theorem gives
a closed-form estimate.  Setting the kinetic and potential energy scales
equal yields
\begin{equation}\label{eq:opt_alpha_mono}
  \alpha_{\mathrm{opt}}^{(\mathrm{mono})}
  = \left[\frac{N c_N\,(N-1)!!}{2\left(2\mu\right)^{N/2}}\right]^{2/(N+2)},
\end{equation}
where $(N-1)!! = 1\cdot 3\cdot 5\cdots(N-1)$ is the double factorial.
For $N=4$ (quartic double-well at the barrier) this recovers the standard
result $\alpha_{\mathrm{opt}}=(3c_4/\mu)^{1/3}$ up to numerical factors.

\subsubsection{Practical recommendation}

For production runs, fix $\alpha$ by the local harmonic frequency at the
minimum of $V(x)$.  If the potential has minima at $x_{\pm}$, set
\begin{equation}\label{eq:alpha_harm}
  \alpha = \mu\,\omega_{\min},
  \qquad
  \omega_{\min} = \sqrt{\frac{V''(x_+)}{\mu}}.
\end{equation}
For the symmetric double-well this reduces to $\alpha=\mu\omega_0$ with
$\omega_0 = \sqrt{4b}$ (see Section~\ref{sec:potential}).

% ============================================================
\subsection{Parity decomposition for symmetric potentials}\label{sec:parity_gen}
% ============================================================

When all odd-index coefficients vanish ($c_1=c_3=\cdots=0$) the potential is
even: $V(-x)=V(x)$.  The Hamiltonian then commutes with the parity operator
$\hat\Pi$, and the eigenstates are either even ($\phi_n$ with $n$ even) or
odd ($\phi_n$ with $n$ odd).  The Hamiltonian matrix block-diagonalises into
two independent $(M/2)\times(M/2)$ sub-matrices:
\begin{equation}
  \mathbf{H}
  = \mathbf{H}^{(+)} \oplus \mathbf{H}^{(-)},
\end{equation}
where $\mathbf{H}^{(+)}$ acts on even-index basis functions and
$\mathbf{H}^{(-)}$ on odd-index ones.  Exploiting this structure reduces the
diagonalisation cost by a factor of $\sim 4$.

% ============================================================
\subsection{Number of basis functions and convergence}\label{sec:conv_gen}
% ============================================================

The number of basis functions $M$ required for a given accuracy depends on
two competing factors:
\begin{enumerate}
  \item \textbf{Energy range:} All HO basis functions with energies
        $E_n^{(\mathrm{HO})}\lesssim E_{\mathrm{max}}$ should be included,
        where $E_{\mathrm{max}}$ is the highest eigenvalue of interest.
  \item \textbf{Coupling range:} The bandwidth $N$ means each function couples
        to its $N$-th nearest neighbours; including too few functions truncates
        off-diagonal couplings.
\end{enumerate}

A practical criterion: choose $M$ such that
\begin{equation}\label{eq:M_criterion}
  E_{M-1}^{(\mathrm{HO})} = \alpha\hb(M-\tfrac{1}{2}) \gg V_{\mathrm{max}},
\end{equation}
where $V_{\mathrm{max}}$ is the energy of the highest eigenstate to be
resolved.  Convergence should always be verified by increasing $M$ until the
eigenvalues of interest change by less than the desired tolerance.

% ============================================================
\subsection{General wavepacket construction}\label{sec:gen_wp}
% ============================================================

The general initial state $\ket{\Psi(0)}$ is specified as a superposition of
the first $M$ eigenstates $\{\ket{\psi_j}\}_{j=0}^{M-1}$:
\begin{equation}\label{eq:gen_wp}
  \ket{\Psi(0)} = \sum_{j=0}^{M-1} c_j\,\ket{\psi_j},
  \qquad
  \sum_{j=0}^{M-1}|c_j|^2 = 1.
\end{equation}
Three natural construction modes are supported by \texttt{QuTu}.

\subsubsection{Mode~1: localized Gaussian}\label{sec:wp_gaussian}

Centered at position $x_0$ with momentum $p_0$ and width $\sigma$:
\begin{equation}\label{eq:wp_gauss}
  \Psi(x,0) = \mathcal{N}\exp\!\left(-\frac{(x-x_0)^2}{4\sigma^2}
              + i p_0 x\right), \qquad
  \mathcal{N} = (2\pi\sigma^2)^{-1/4}.
\end{equation}
The expansion coefficients are
\begin{equation}\label{eq:cj_gauss}
  c_j = \braket{\psi_j}{\Psi(0)}
      = \int_{-\infty}^{\infty} \psi_j^*(x)\,\Psi(x,0)\,dx,
\end{equation}
evaluated numerically on a grid fine enough to resolve both $\psi_j$ and
$\Psi(x,0)$.

\subsubsection{Mode~2: energy-filtered state}\label{sec:wp_energy}

Populate all eigenstates in an energy window
$[E_{\mathrm{low}},E_{\mathrm{high}}]$ with equal weight:
\begin{equation}\label{eq:wp_energy}
  c_j =
  \begin{cases}
    \mathcal{N}^{-1} & \text{if } E_{\mathrm{low}} \leq E_j \leq E_{\mathrm{high}},\\
    0 & \text{otherwise,}
  \end{cases}
  \quad
  \mathcal{N}^{2} = \#\{j : E_{\mathrm{low}}\leq E_j\leq E_{\mathrm{high}}\}.
\end{equation}
This creates a \emph{quasi-classical} packet that samples a controlled energy
range without presupposing spatial localisation.

\subsubsection{Mode~3: user-defined amplitudes}\label{sec:wp_user}

The user specifies the complex amplitudes $\{c_j\}$ directly.  The code
renormalises: $c_j \leftarrow c_j/\|\mathbf{c}\|$.  This mode is useful for
initialising the two-state packet $c_0=c_1=1/\sqrt{2}$ (symmetric
tunneling) or the four-state packet of Section~\ref{sec:wavepackets}.

% ============================================================
\subsection{Time evolution and observables}\label{sec:gen_time_evo}
% ============================================================

Given the general expansion~\eqref{eq:gen_wp}, the time-evolved state is
\begin{equation}\label{eq:gen_psi_t}
  \ket{\Psi(t)} = \sum_{j=0}^{M-1} c_j\,e^{-iE_j t/\hb}\,\ket{\psi_j}.
\end{equation}
All observables of Section~\ref{sec:observables} carry over without
modification: survival probability, position expectation value, tunneling
splitting, and turning points are each determined by the same formulas with
$\{E_j,\psi_j,c_j\}$ now drawn from the general diagonalisation.

\begin{remark}
  The tunneling splitting $\Delta E=E_1-E_0$ and half-period
  $T_{\mathrm{tun}}=h/\Delta E$ remain meaningful for any symmetric double-well
  regardless of the polynomial degree $N$, as long as $E_0$ and $E_1$ are
  below the barrier.
\end{remark}

% ============================================================
\subsection{Implementation summary}\label{sec:gen_impl}
% ============================================================

Table~\ref{tab:gen_impl} summarises the algorithmic flow for the general
polynomial case as implemented in \texttt{QuTu}.

\begin{table}[htbp]
  \centering
  \caption{%
    Step-by-step algorithm for the general polynomial potential solver.
    All matrix operations are $O(M^2)$ in memory; the dense \texttt{dsyev}
    eigensolver is $O(M^3)$.  A banded \texttt{dsbev} path is planned.%
  }
  \label{tab:gen_impl}
  \ra{1.3}
  \begin{tabular}{clp{7cm}}
    \toprule
    Step & Operation & Notes \\
    \midrule
    1 & Read $\{c_k\}$, $N$, $M$, $\mu$, $\alpha$ &
        User input; $\alpha$ from Eq.~\eqref{eq:alpha_harm} if not given \\
    2 & Compute $X^{(k)}$ for $k=0,\ldots,N$ &
        Master recursion~\eqref{eq:master_rec}, $O(kM^2)$ per $k$ \\
    3 & Assemble $V_{mn} = \sum_k c_k X^{(k)}_{mn}$ &
        Linear combination, $O(NM^2)$ \\
    4 & Add $T_{mn}$; form $H_{mn}$ &
        Analytic kinetic matrix, Section~\ref{sec:HO} \\
    5 & Diagonalise $\mathbf{H}$ &
        LAPACK \texttt{dsyev} (symmetric dense); banded \texttt{dsbev} planned \\
    6 & Construct initial state $\ket{\Psi(0)}$ &
        One of the three modes, Section~\ref{sec:gen_wp} \\
    7 & Propagate: $c_j(t)=c_j e^{-iE_j t/\hb}$ &
        Analytic; $O(M)$ per time step \\
    8 & Compute observables &
        Section~\ref{sec:observables} and~\ref{sec:additional_obs} \\
    \bottomrule
  \end{tabular}
\end{table}
